\section{Background and Related Work}
\label{sec:background}

\subsection{Markdown and LLM Training Data}

Markdown was published in 2004 by John Gruber, in collaboration with Aaron Swartz, as ``a text-to-HTML conversion tool for web writers''~\cite{gruber2004markdown}, formalizing conventions that writers were already employing informally in plain text. It has since become the dominant formatting layer for structured text on the internet, adopted by GitHub~\cite{github2024octoverse}, Stack Overflow, Reddit, Discord, Slack, and numerous documentation platforms. Two major standardization efforts---CommonMark~\cite{commonmark2021} and GitHub Flavored Markdown~\cite{gfm2024}---have formalized the syntax, while Pandoc~\cite{pandoc2024} provides the most comprehensive processing pipeline, including smart typography extensions that convert certain character sequences (notably \texttt{---} to em dash) during rendering.

The training data for large language models draws heavily from this markdown-saturated web. The Pile~\cite{gao2020pile}, an 800GB dataset widely used for language model training, includes GitHub repositories (where essentially all prose content is markdown), Stack Exchange (markdown-formatted), and a broad web crawl in which markdown-formatted content is disproportionately represented among high-quality, well-structured text. RedPajama~\cite{together2023redpajama} and RefinedWeb~\cite{penedo2023refinedweb} follow similar patterns. Audits of large web corpora~\cite{dodge2021documenting, longpre2023pretrainers} confirm that high-quality text sources are disproportionately drawn from platforms where markdown is the native formatting layer.

Precise estimates of the proportion of markdown-formatted text in LLM training corpora are not publicly available. However, the structural evidence is suggestive: the platforms that generate the highest-quality, most heavily structured text online---developer documentation, technical writing, academic and educational content---overwhelmingly use markdown. Training data curation processes that select for text quality therefore select, implicitly, for markdown-formatted text. The structural conventions of markdown---headings for hierarchical organization, dashes for boundaries and separators, asterisks for emphasis, lists for discrete enumeration---are not peripheral features of the training distribution. They are the organizational grammar of the highest-quality text in the corpus.

\subsection{RLHF and Output Style}

Reinforcement learning from human feedback (RLHF)~\cite{ouyang2022training, christiano2017deep} is one of several post-training mechanisms---alongside supervised fine-tuning~\cite{wei2022finetuned, sanh2022multitask} and constitutional methods~\cite{bai2022training}---by which base language models are adapted to produce outputs that users find helpful, clear, and well-organized. In the RLHF pipeline, human evaluators rate model outputs, and these ratings are used to train a reward model that guides further model optimization~\cite{bai2022training, zheng2024judging}.

The stylistic consequences of RLHF are well-documented but incompletely understood. RLHF-tuned models produce outputs that are more structured, more organized, and more likely to include formatting conventions (headers, lists, bold emphasis) than their base model counterparts. The mechanism is straightforward: human evaluators, who are disproportionately drawn from technical and educated populations, tend to rate well-organized, clearly structured prose more highly. RLHF therefore selects for the structural register---the ``markdown voice''---that these evaluators prefer.

This has direct implications for em dash frequency. If evaluators reward prose that reads as precise, articulate, and structurally aware---and if em-dash-heavy prose is perceived as having these qualities---then RLHF will amplify em dash use even in models whose base weights do not predispose them to it. The distinction between base model tendencies and RLHF-amplified tendencies is central to the genealogy we propose in Section~\ref{sec:genealogy}.

\subsection{Prior Work on AI Stylistic Artifacts}

Academic work on the stylistic properties of AI-generated text has focused primarily on detection. DetectGPT~\cite{mitchell2023detectgpt} uses probability curvature to identify machine-generated text without requiring access to the generating model. Kirchenbauer et al.~\cite{kirchenbauer2023watermark} propose a watermarking scheme that embeds detectable statistical patterns in LLM output. Ippolito et al.~\cite{ippolito2020automatic} find that automatic detection is easiest when generated text is most convincing to humans, while Liang et al.~\cite{liang2023gpt} demonstrate that GPT detectors exhibit systematic biases, raising questions about the reliability of distributional approaches. These methods operate at the distributional level---they detect machine-generated text through statistical fingerprints rather than through analysis of specific stylistic features.

To our knowledge, no prior work has proposed a \textit{mechanistic} account of why specific punctuation patterns (such as em dash overuse) arise in LLM output, connecting them to the structural conventions of training data. The observation that ``AI overuses em dashes'' is widely noted in informal discourse but has not been the subject of formal analysis. Similarly, the observation that ``AI thinks in markdown'' is widely discussed in developer communities but has not been connected to prose-level consequences. This paper bridges these two informal discourses and proposes the first mechanistic account linking training data format to specific punctuation behavior in generated prose.
